\documentclass[10pt,a4paper]{amsart}
\usepackage[margin=19mm]{geometry}
\usepackage{amsmath,amssymb,mathtools}
\usepackage[scr=rsfs,cal=euler]{mathalfa}
\usepackage{xfrac}
\usepackage[unicode,hidelinks,bookmarks=false]{hyperref}
\newcommand{\ie}{i.e.\ }
\newcommand{\cf}{cf.\ }
\newcommand{\resp}{respectively\ }
\newcommand{\Subsection}{\subsection}
\providecommand{\nobreakdash}{\nobreak\hbox{-}}
\setlength{\emergencystretch}{2em}
\multlinegap0pt
\usepackage{bm}
\usepackage{bbm}
\usepackage{dsfont}
\usepackage{xspace}
\usepackage{cancel}
\usepackage{mathtools}
\usepackage{amsthm}
\makeatletter
\@ifundefined{thm}{\newtheorem{thm}{Theorem}}{}
\@ifundefined{lem}{\newtheorem{lem}[thm]{Lemma}}{}
\makeatother
\DeclareMathOperator{\supp}{supp}
\newcommand{\Ga}{\Gamma}
\newcommand{\Tree}{\mathrm{T}}
\newcommand{\acts}{\curvearrowright}
\newcommand{\cU}{\mathcal{U}}
\newcommand{\ip}[1]{\mathopen{\langle}#1\mathclose{\rangle}}
\newcommand{\vp}{\varphi}
\title[Proximality and selflessness for group C*-algebras]{Proximality and selflessness for group C*-algebras}
\author{Narutaka Ozawa}
\date{Source: arXiv:2508.07938v8 (CC BY 4.0)}
\begin{document}
\maketitle

\noindent\emph{Source and licence.} Proximality and selflessness for group C*-algebras. Version arXiv:2508.07938v8 (CC BY 4.0), distributed under the Creative Commons Attribution 4.0 International licence, as recorded in the arXiv licence metadata for this exact version and shown on the abstract page https://arxiv.org/abs/2508.07938v8. Full rights notice: licenses/arxiv-2508.07938-CC-BY-4.0.txt.\par\smallskip

\noindent\textit{Editorial scope.} Counted unit: the Thm whose source label is \texttt{thm:space} in the article (source0.tex lines 578--626, source label \texttt{thm:space}), delivered together with the article's own complete original proof of it (source0.tex lines 628--749, one \texttt{proof} environment of 122 lines). No mathematics is written, repaired or re-derived: statement and proof are the article's own text line for line. Required context retained uncounted: lem:labelequiv (lines 465--471); lem:pingpong (lines 482--488). Every definition, display and label of those lines is retained unchanged. Omitted from this package and not redistributed: 1--464, 472--481, 489--577, 627--627, 750--1656. Nothing inside a retained excerpt is omitted or altered: every retained body is delivered exactly as the article's lines are written, so no omission receipt of any kind accompanies this package. Citations: the retained excerpts carry no citation command at all, so no bibliography entry of the article is transcribed and no reference is redistributed. Reference adaptation: none was needed and none was made; every reference inside the delivered unit resolves inside it, so no unresolved reference or citation is produced. Printed slips kept literal: no printed word, symbol, hypothesis or number of the counted statement or of its proof was altered. Layout and class: the article's own class is \texttt{amsart} with packages including amssymb, mathtools, accents; the wrapper is the lane's pinned \texttt{amsart} editorial wrapper at 10pt on A4, which restates the theorem-like environments the retained text needs and adds the article's own macro definitions that the retained text needs (\texttt{\textbackslash Ga}, \texttt{\textbackslash Tree}, \texttt{\textbackslash acts}, \texttt{\textbackslash cU}, \texttt{\textbackslash ip}, \texttt{\textbackslash supp}, \texttt{\textbackslash vp}), with extra wrapper packages (bm, bbm, dsfont, xspace, cancel, mathtools). The delivered numbering is the unit's own. Assets: no retained span contains a figure environment, a graphics-inclusion command, a PDF-page inclusion, a table or a TikZ picture; no image file is copied into this package, the package declares no asset dependency and no image file is redistributed with it. Licence: version arXiv:2508.07938v8 of this article is distributed under the Creative Commons Attribution 4.0 International licence, as recorded in the arXiv licence metadata for this exact version and shown on the abstract page https://arxiv.org/abs/2508.07938v8. A full rights notice is delivered at licenses/arxiv-2508.07938-CC-BY-4.0.txt. Characters: every retained excerpt is pure ASCII, so the delivered engine, XeLaTeX with its default Latin Modern text font, reports no missing character.
\begin{lem}\label{lem:labelequiv}
Let $(t,s)\in (\Tree\times\bar{\Tree})$ 
be a non-diagonal point. 
For $a\in\Ga$, one has $\ell(o,as)=a\ell(o,s)$ 
and $\ell(at,as) = \ell(t,s)$ for $t\neq o$. 
Also, $\ell(zt,zs)=\ell(t,s)$. 
\end{lem}

\begin{lem}\label{lem:pingpong}
Let $g=a_1z^{\epsilon_1} \cdots z^{\epsilon_l}a_{l+1}$ 
be an element in the reduced form with $l\geq1$ 
and $x\in X \setminus\{a_{l+1}^{-1} z_{\bar{\epsilon}_l}\}$, 
where $\bar{\epsilon}$ is the opposite of $\epsilon$. 
Then $\vp_n(g) x \to a_1z_{\epsilon_1}$ as $n\to\infty$.
\end{lem}

\begin{thm}\label{thm:space}
In the above setting for $\tilde{\Ga}=\Ga*\ip{z}\acts\Tree$, the following hold.
\begin{enumerate}
\item\label{p1}
$X_\Tree = \theta(\Tree,X) \cup \theta(\partial\Tree)$.
\item\label{p2}
$\theta(\Tree,X) \cap \theta(\partial\Tree) = \emptyset$.
\item\label{p3}
For $s\neq t$ and $x,y\in X$, 
one has $\theta(s,x)=\theta(t,y)$ if and only if 
$s$ and $t$ are adjacent and $x=\ell(s,t)$ and $y=\ell(t,s)$. 
\item\label{p4}
For $a\in\Ga$, one has 
$a\theta(o,x)=\theta(o,ax)$ and $a\theta(s,x)=\theta(as,x)$ 
for $s\neq o$.
\item\label{p5}
$z\theta(s,x)=\theta(zs,x)$. 
\item\label{p6}
$\theta\colon \partial\Ga\to X_\Tree$ is $\tilde{\Ga}$-equivariant.
\item\label{p7}
A sequence $(\theta(t_n,x_n))_n$ 
converges to $\theta(s,x)$ in $X_\Tree$ if and only if 
either 
$(t_n)_n$ converges to $s$ in $\bar{\Tree}$ and 
the sequence $(y_n)_n$, defined by 
$y_n= x_n$ if $t_n=s$ and $y_n = \ell(s,t_n)$ if $t_n\neq s$,
converges to $x$ in $X$; 
or there is an adjacent point $s_1$ to $s_0 \coloneq s$ in $\Tree$ 
such that $t_n\in\{s_0,s_1\}$ eventually, $x=\ell(s_0,s_1)$, 
and the subsequences $(x_n)_{\{n : t_n=s_i\}}$ 
(possibly there is only one of them) converge 
to $\ell(s_i,s_j)$ for $\{i,j\}=\{0,1\}$.  
A sequence $(\theta(t_n,x_n))_n$ 
converges to $\theta(\omega)$ in $X_\Tree$ if and only if 
$(t_n)_n$ converges to $\omega$ in $\bar{\Tree}$.
\item\label{p8}
A sequence $\theta(\omega_n)$ 
converges to $\theta(s,x)$ in $X_\Tree$ if and only if 
$(\omega_n)_n$ converges to $s$ in $\bar{\Tree}$ and 
and $(\ell(s,\omega_n))_n$ converges to $x$ in $X$.
\item\label{p9}
One has 
$\theta(\Tree,\Ga\{z_\pm\})=\tilde{\Ga}\theta(o,\{z_\pm\})$. 
The map 
$\beta\colon X_\Tree \setminus \tilde{\Ga}\theta(o,\{z_\pm\}) 
\to\bar{\Tree}$, given by $\beta(\theta(s,x))=s$ and $\beta(\theta(\omega))=\omega$, is 
$\tilde{\Ga}$-equivariant and continuous. 
\end{enumerate}
\end{thm}

\begin{proof}
We prove (\ref{p3}). The proof of $(\ref{p2})$ is similar. 
Suppose that $s$ and $t$ are adjacent 
and $x=\ell(s,t)$ and $y=\ell(t,s)$. 
Then $\theta(s,x)_s=x=\ell(s,t)=\theta(t,y)_s$, 
$\theta(s,x)_t=\ell(t,s)=y=\theta(t,y)_t$, and 
$\theta(s,x)_r = \theta(r,s)=\theta(r,t)=\theta(t,y)_r$ 
for $r\in\Tree\setminus\{s,t\}$, because $s$ and $t$ 
are on the same side of $r$. 
This shows $\theta(s,x)=\theta(t,y)$. 
Conversely, let $(s,x)$ and $(t,y)$ be given. 
Suppose that there is $r\in\Tree$ strictly between $s$ and $t$. 
Then $\theta(s,x)\neq\theta(t,y)$ since 
$\theta(s,x)_r = \ell(r,s) \neq \ell(r,t) = \theta(t,y)_r$. 
Suppose next that $s$ and $t$ are adjacent. 
Then $\theta(s,x)_s=x$, $\theta(s,x)_t=\ell(t,s)$ 
and $\theta(t,y)_s=\ell(s,t)$, $\theta(t,y)_t=y$. 
Hence $\theta(s,x)=\theta(t,y)$ only if $x=\ell(s,t)$ and $y=\ell(t,s)$.
This proves (\ref{p3}). 

The assertions (\ref{p4}-\ref{p6}) 
follow from Lemma~\ref{lem:labelequiv}; 
E.g., for $a\in\Ga$ and $s\in\Tree\setminus\{ o\}$, one has
\[
(a\theta(s,x))_o=a\theta(s,x)_o=a\ell(o,s)=\ell(o,as)=\theta(as,x)_o
\]
and for $t\neq o$ 
\[
(a\theta(s,x))_t = \theta(s,x)_{a^{-1}t}
 = \left\{\begin{array}{ll}
     x & \mbox{ if $t=as$} \\
     \ell(a^{-1}t,s) = \ell(t,as) & \mbox{ if $t\neq as$} 
   \end{array}\right\}
 =\theta(as,x)_t.
\]

The ``if'' part of (\ref{p7}-\ref{p8}) are easy. 
For the ``only if'' part, suppose that 
$(\theta(t_n,x_n))_n$ converges to $\theta(s,x)$. 
If $(t_n)_n$ is converges to some point $s'$, 
then by the ``if'' part and the compactness of $X$, 
one has $\theta(t_n,x_n) \to \theta(s',y)$, 
where $y$ is the limit of $y_n \coloneq x_n$ if $t_n=s'$ and 
$y_n \coloneq \ell(s',t_n)$ if $t_n\neq s'$. 
Since $\theta(s',y)=\theta(s,x)$, one has by (\ref{p3}) 
either $s'=s$ or $s'$ is an adjacent point to $s$ 
that satisfies $x=\ell(s,s')$ and $y=\ell(s',s)$. 
This proves the ``only if'' part for the case $(t_n)_n$ is convergent. 
Now consider the case where $(t_n)_n$ is not convergent. 
By the previous case, the only limit points are 
either $s$ or the unique adjacent point $s'$ 
that satisfies $x=\ell(s,s')$. 
The other claimed conditions also follow. 
The rest of the proof of (\ref{p7}-\ref{p8}) is similar. 

We prove (\ref{p1}). 
We first prove that $\theta(\Tree,X)\subset X_\Tree$. 
Let $s\in\Tree$ and $x\in X\setminus\Ga\{z_\pm\}$ be given.
We consider the character $\chi$ on 
$C(X_\Tree)\subset C(X)^\cU$ 
arising from the sequence $(\vp_n(\rho(s))x)_n$. 
Then by Lemma~\ref{lem:pingpong} one has 
\[
\chi(f^{(t)}) = \lim\nolimits_{\cU} \sigma_{\vp_n(\rho(t))}(f)(\vp_n(\rho(s))x) 
 = \lim\nolimits_{\cU} f(\vp_n(\rho(t)^{-1}\rho(s))x) 
 = f(\ell(t,s))
\]
for $t\neq s$ and $\chi(f^{(s)})=f(x)$. 
This means $\chi$ corresponds to 
the point $\theta(s,x)$ in $X_\Tree$. 
Since $\theta(s,\,\cdot\,)\colon X\to X^\Tree$ is 
continuous and $X\setminus\Ga\{z_\pm\}$ is dense in $X$, 
one sees that $\theta(s,X)\subset X_\Tree$. 
Similarly for $\theta(\partial\Tree)$. 
It follows that $\theta(\Tree,X)\cup\theta(\partial\Tree)\subset X_\Tree$. 

It remains to prove the opposite inclusion. 
Let $\mathbf{x} \coloneq (x_t)_t\in X^\Tree$ be given and assume that 
$\mathbf{x} \in X_\Tree$. 
We define $w\colon \Tree \to \Tree$, 
``the orientation of $\mathbf{x}$ at $t$'', as follows. 
If $x_t \in X\setminus\Ga\{z_\pm\}$, then $w(t) \coloneq t$. 
If $x_t \in \Ga\{ z_\pm\}$, then we define $w(t)$ to be the unique 
point that is adjacent to $t$ and satisfies $\ell(t,w(t))=x_t$. 
We claim that if $s,t\in\Tree$ are adjacent, then 
it cannot happen that $w(s)\neq t$ and $w(t)\neq s$ simultaneously. 
Suppose for a contradiction that it happened. 
We may assume that $s$ is between $o$ and $t$ 
and $\rho(t)=\rho(s) a z$ with no cancellation. 
Thus one has $x_s \neq az_+$ and $x_t \neq z_-$. 
Let $f\in C(X)$ be such that 
$f(x_s)=1$ and $f=0$ on a neighborhood of $az_+$. 
Similarly let $g$ be such that $g(x_t)=1$ and $g=0$ 
on a neighborhood of $z_-$. 
We consider the function $f^{(s)}g^{(t)}$ on $X^\Tree$. 
It satisfies $(f^{(s)}g^{(t)})( \mathbf{x} ) = f(x_s)g(x_t) = 1$.
However, 
\[
(f^{(s)}g^{(t)})|_{X_\Tree} = Q(f^{(s)}g^{(t)}) 
 = \sigma_{\rho(s)}(f)\sigma_{\rho(t)}(g)
 = \sigma_{\rho(s)}( f \sigma_{az}(g) ) = 0
\]
since $(\supp f)  \cap az_n(\supp g) = \emptyset$ eventually. 
This means that $\mathbf{x} \in X^\Tree\setminus X_\Tree$, 
a contradiction. 
Thus we have proved the claim for adjacent 
points $s,t\in\Tree$ that $w(s)\neq t$ implies $w(t)=s$. 
This yields three possibilities. 
(1): There is $s\in \Tree$ such that $w$ is oriented to 
the $s$ direction. 
(2): There is $\omega\in\partial\Tree$ such that 
$w$ is oriented to the $\omega$ direction. 
(3): There is a pair $s_1,s_2$ of adjacent points such that 
$w$ is oriented to $\{s_1,s_2\}$ direction 
and $w(s_i)=s_j$ for $\{ i, j \}=\{1,2\}$. 
In the first case, one has $\mathbf{x}=\theta(s,x_s)$. 
In the second case, one has $\mathbf{x}=\theta(\omega)$. 
In the third case, one has $\mathbf{x}=\theta(s_1,\ell(s_1,s_2))=\theta(s_2,\ell(s_2,s_1))$.
It follows that $\theta(\Tree,X)\cup\theta(\partial\Tree)\supset X_\Tree$. 

Finally, the assertion (\ref{p9}) is a consequence of (\ref{p1}-\ref{p8}).
\end{proof}

\end{document}
